% Week 6 Session B Lab 4 — Flood Inundation Analysis
% Build: pdflatex Week6_SessionB_Lab4_Report.tex from Lab-reports/ (twice)
% Assets:
%   week6_session_b_lab4_terminal.png  (all four commands, one screenshot)
%   week6_session_b_lab4_dem.png
%   week6_session_b_lab4_comparison.png
%   week6_session_b_lab4_monotonic.png
% Code: week6_session_b_lab4_files/ (ASCII-safe for VerbatimInput)
%
\documentclass[11pt,a4paper]{article}

\usepackage[margin=2.5cm]{geometry}
\usepackage{graphicx}
\newcommand{\termfig}[1]{%
  \includegraphics[width=\linewidth,height=0.55\textheight,keepaspectratio]{#1}%
}
\newcommand{\plotfig}[1]{%
  \includegraphics[width=\linewidth,height=0.42\textheight,keepaspectratio]{#1}%
}
\newcommand{\plotfigwide}[1]{%
  \includegraphics[width=\linewidth,height=0.38\textheight,keepaspectratio]{#1}%
}
\usepackage{booktabs}
\usepackage{xurl}
\usepackage{hyperref}
\usepackage{parskip}
\usepackage{enumitem}
\usepackage{xcolor}
\usepackage{fvextra}
\usepackage[most]{tcolorbox}

\newcommand{\studentname}{Mahmudul Hasan}
\newcommand{\studentid}{4125999049}
\newcommand{\reportdate}{May 21, 2026}

\makeatletter
\g@addto@macro\UrlBreaks{\do\ }
\makeatother

\newcommand{\LabCodeRoot}{week6_session_b_lab4_files}

\newcommand{\filebox}[2]{%
\begin{tcolorbox}[
  enhanced,
  colback=gray!6,
  colframe=black!55,
  title={#1},
  fonttitle=\bfseries\ttfamily\footnotesize,
  arc=1.5mm,
  boxrule=0.7pt,
  breakable,
  width=\linewidth,
  before skip=0.8em,
  after skip=0.8em,
]
\VerbatimInput[fontsize=\footnotesize,breaklines=true,breakanywhere=true]{\LabCodeRoot/#2}
\end{tcolorbox}%
}

\title{Lab Report: Week 6 Session B Lab 4\\Flood Inundation Analysis}
\author{\studentname \quad (Student ID: \studentid)}
\date{\reportdate}

\begin{document}
\maketitle

\begin{abstract}
This report documents Week~6 Session~B Lab~4: \textbf{synthetic DEM} generation, \textbf{flat-water inundation} mapping, depth statistics, and a \textbf{rising-stage} sensitivity over flood levels 40--60~m. A 100$\times$100 grid with 1~m cells and random seed 42 yields elevations 35.96--68.81~m. Wet cells use \texttt{elevation <= flood\_level}; depth is \texttt{max(0, flood\_level - elevation)} on wet cells only. At 40, 50, and 60~m, flooded area is 2.41\%, 25.25\%, and 82.93\% with mean wet depth 1.18, 4.86, and 8.33~m. Exercise~4 swept levels in 1~m steps: \textbf{Monotonicity check: PASS}. Evidence: one combined terminal screenshot (Figure~\ref{fig:terminal}), three figures, and code in \texttt{week6\_session\_b\_lab4/} (Appendix~\ref{app:core}--\ref{app:promptlog}). OpenCode was not required.
\end{abstract}

\section{Introduction}
Flood risk screening often starts from a digital elevation model (DEM) and a assumed water surface elevation. This lab practices GIS-style inundation on a \textbf{synthetic terrain}: generate DEM data, apply a simple depth rule, visualize extent, and verify that flooded area grows monotonically as stage rises---after reservoir optimization (Week~6 Session~A Lab~3) and hydrologic modeling (Weeks~3--5).

\section{Objectives}
\begin{itemize}[noitemsep]
  \item Process (synthetic) DEM data for spatial analysis (Exercise~1).
  \item Implement \texttt{simulate\_flood} and area/depth statistics (Exercise~2).
  \item Create publication-style inundation maps (Exercise~3).
  \item Simulate rising water levels and check monotonic flooded area (Exercise~4).
\end{itemize}

\section{Methods}

\subsection{Flood level definition}
Following the lab guide, \texttt{flood\_level} is \textbf{absolute stage} (m) in the same vertical datum as the DEM---not depth above a channel thalweg unless redefined.

\subsection{Inundation rules}
\begin{itemize}[noitemsep]
  \item \textbf{Wet mask:} \texttt{dem <= flood\_level} (inclusive at the water surface).
  \item \textbf{Depth:} \texttt{max(0, flood\_level - dem)} on wet cells; \texttt{0} elsewhere.
  \item \textbf{Grid:} 100$\times$100 cells, \textbf{1~m} spacing; seed \texttt{42} for reproducibility.
\end{itemize}

\subsection{Scripts}
\begin{table}[htbp]
  \centering
  \caption{Submitted files in \texttt{week6\_session\_b\_lab4/}.}
  \label{tab:files}
  \small
  \begin{tabular}{@{}ll@{}}
    \toprule
    File & Role \\
    \midrule
    \texttt{flood\_inundation.py} & Core \texttt{simulate\_flood}, stats, monotonicity check \\
    \texttt{generate\_dem.py} & Exercise~1: DEM + heatmap \\
    \texttt{run\_flood\_tests.py} & Exercise~2: levels 40, 50, 60~m \\
    \texttt{visualize\_floods.py} & Exercise~3: side-by-side maps (300 DPI) \\
    \texttt{rising\_water.py} & Exercise~4: 40--60~m sweep + plot \\
    \texttt{data/dem.npy}, \texttt{figures/*.png} & Saved outputs \\
    \texttt{prompt\_log.md} & AI interaction log \\
    \bottomrule
  \end{tabular}
\end{table}

\section{Results}

\subsection{Terminal verification}
Figure~\ref{fig:terminal} shows one screenshot with all four commands (\texttt{generate\_dem.py}, \texttt{run\_flood\_tests.py}, \texttt{visualize\_floods.py}, \texttt{rising\_water.py}), matching the lab guide's optional single-screenshot submission style.

\begin{figure}[htbp]
  \centering
  \termfig{week6_session_b_lab4_terminal.png}
  \caption{Terminal evidence: DEM saved (35.96--68.81~m), flood table at 40/50/60~m, comparison figure saved, monotonicity \textbf{PASS}.}
  \label{fig:terminal}
\end{figure}

\subsection{Exercise 2: flood levels}
\begin{table}[htbp]
  \centering
  \caption{Flooded area and mean depth on wet cells (from \texttt{run\_flood\_tests.py}).}
  \label{tab:levels}
  \begin{tabular}{@{}rrrr@{}}
    \toprule
    Flood level (m) & Flooded (\%) & Mean depth (m) & Wet cells \\
    \midrule
    40 & 2.41 & 1.183 & 241 \\
    50 & 25.25 & 4.863 & 2525 \\
    60 & 82.93 & 8.329 & 8293 \\
    \bottomrule
  \end{tabular}
\end{table}

Flooded fraction increases sharply between 50 and 60~m because most valley cells lie below 60~m but many ridges remain dry until stage approaches the DEM maximum.

\subsection{Exercise 1: DEM}
\begin{figure}[htbp]
  \centering
  \plotfig{week6_session_b_lab4_dem.png}
  \caption{Synthetic DEM heatmap (\texttt{figures/dem\_heatmap.png}).}
  \label{fig:dem}
\end{figure}

If \texttt{week6\_session\_b\_lab4\_dem.png} is missing on Overleaf, use \texttt{week6\_session\_b\_lab4\_files/figures/dem\_heatmap.png}.

\subsection{Exercise 3: inundation maps}
\begin{figure}[htbp]
  \centering
  \plotfigwide{week6_session_b_lab4_comparison.png}
  \caption{Side-by-side inundation at 40, 50, and 60~m (grayscale DEM, blue depth overlay).}
  \label{fig:comparison}
\end{figure}

\subsection{Exercise 4: rising water}
\begin{figure}[htbp]
  \centering
  \plotfig{week6_session_b_lab4_monotonic.png}
  \caption{Flooded domain percentage vs.\ flood level (40--60~m, step 1~m).}
  \label{fig:monotonic}
\end{figure}

At 40~m flooded area is 2.41\%; at 60~m it is 82.93\%. \texttt{check\_monotonic\_area} reports \textbf{PASS} for the 1~m step series.

\section{Analysis and validation}

\textbf{Inclusive wet mask:} Using \texttt{<=} matches a flat water surface touching cell centers at exactly stage elevation; switching to \texttt{<} would shrink area slightly---document the choice in \texttt{prompt\_log.md}.

\textbf{Monotonicity:} For a fixed DEM and non-decreasing \texttt{flood\_level}, the wet-cell count cannot shrink; PASS confirms implementation consistency.

\textbf{Mean depth:} Averaged over \textbf{wet cells only}, so it rises with stage both because deeper inundation and because lower-elevation cells join the wet set.

\textbf{Visualization warning:} \texttt{visualize\_floods.py} may emit a harmless \texttt{tight\_layout} warning with three subplots; the 300~DPI PNG still saves correctly (\texttt{bbox\_inches=tight}).

\section{Discussion}
\begin{enumerate}[noitemsep]
  \item A single terminal screenshot is sufficient when figures carry the visual evidence (lab guide: screenshot optional).
  \item Synthetic DEM with valleys and hills makes inundation statistics interpretable without downloading real GIS tiles.
  \item The 50--60~m jump in flooded percent shows sensitivity to vertical datum choice in real projects.
  \item No OpenCode required; generic AI prompts from the lab guide suffice for scaffolding.
\end{enumerate}

\section{Problems encountered}
Only a matplotlib \texttt{tight\_layout} user warning on the three-panel figure; output file verified on disk. No API keys or external DEM download needed.

\section{Lessons learned}
Document \texttt{flood\_level} definition and wet/dry convention before coding. Fix the random seed early. Use one terminal screenshot plus saved PNGs for submission. If monotonicity fails, inspect mask logic and level ordering before adjusting terrain.

\section{Conclusion}
Week~6 Session~B Lab~4 objectives were met: synthetic DEM, inundation logic, comparison maps, and monotonic rising-stage analysis with documented results in \texttt{prompt\_log.md}. Submitted: \texttt{week6\_session\_b\_lab4/}, Figure~\ref{fig:terminal}, and Figures~\ref{fig:dem}--\ref{fig:monotonic}.

\appendix

\section{Core logic: \texttt{flood\_inundation.py}}
\label{app:core}
\filebox{\texttt{flood\_inundation.py} (full file)}{flood_inundation.py}

\section{DEM generator: \texttt{generate\_dem.py}}
\label{app:dem}
\filebox{\texttt{generate\_dem.py} (full file)}{generate_dem.py}

\section{Visualization: \texttt{visualize\_floods.py}}
\label{app:viz}
\filebox{\texttt{visualize\_floods.py} (full file)}{visualize_floods.py}

\section{Rising water: \texttt{rising\_water.py}}
\label{app:rising}
\filebox{\texttt{rising\_water.py} (full file)}{rising_water.py}

\section{Submitted prompt log: \texttt{prompt\_log.md}}
\label{app:promptlog}
\filebox{\texttt{prompt\_log.md} (full file)}{prompt_log.md}

\end{document}
